\documentclass[12pt]{article}
\usepackage{sbc-template}
\usepackage{graphicx,url}
\usepackage[utf8]{inputenc}
\sloppy

\title{Artefato SBSeg 2026 --- Apêndice\\
Artigo \#189 --- TopVenues: An Executable Corpus and Research Tool for
Cybersecurity Literature Reviews}

\author{Sidnei Barbieri, Ágney Roth Ferraz, Lourenço Alves Pereira Júnior}

\address{Instituto Tecnológico de Aeronáutica (ITA)\\
  São José dos Campos --- SP --- Brasil
  \email{sidneisb@ita.br}
}

\begin{document}
\maketitle

\begin{resumo}
Revisões de literatura em segurança dependem de um denominador estável: o
conjunto de artigos elegíveis antes de a triagem começar. Na prática, esse
denominador é remontado a partir de portais, APIs e exportações que mudam ao
longo do tempo, o que torna o corpus resultante difícil de auditar ou reutilizar.
O TopVenues transforma a construção do corpus em um artefato de pesquisa
executável: dado um escopo declarado de veículos e anos, normaliza metadados do
DBLP, enriquece registros com resumos e entradas BibTeX, e materializa um
snapshot SQLite monotônico exposto por linha de comando, aplicação web, busca
ranqueada e exportações orientadas a revisão.
\end{resumo}

\section{Selos Considerados}

Os selos considerados são: \textbf{Disponíveis (SeloD)}, \textbf{Funcionais
(SeloF)}, \textbf{Sustentáveis (SeloS)} e \textbf{Experimentos Reprodutíveis
(SeloR)}.

\section{Informações básicas}

\textbf{Acesso ao artefato.} O artefato é público, sem credenciais, em
repositório estável no GitHub:

\begin{center}
\url{https://github.com/sidneibarbieri/topvenues-tool}
\end{center}

\noindent A versão avaliada é a \texttt{v1.7.0}. O \texttt{README.md} do
repositório é a documentação normativa e segue o modelo mínimo do CTA.

\textbf{Não são necessários recursos externos.} Não há infraestrutura de nuvem,
chave SSH, chave de API, credencial, conta de serviço ou máquina virtual
fornecida pelos autores. O snapshot do corpus acompanha o repositório e a
avaliação é integralmente local.

\textbf{Ambiente usado nos experimentos relatados no artigo.}

\begin{table}[h]
\centering
\begin{tabular}{ll}
\hline
Sistema operacional & macOS 26.5.2 (build 25F84) \\
Kernel & Darwin 25.5.0 \\
Arquitetura & arm64 (Apple Silicon) \\
Processador & Apple M4 Max, 16 núcleos \\
Memória RAM & 64 GB \\
Armazenamento livre & 53 GB \\
Python & 3.14.7 \\
\hline
\end{tabular}
\end{table}

\textbf{Requisitos mínimos para reprodução.} O artefato não exige o hardware
acima. Foram verificados: Linux, macOS ou Windows 10+; Python 3.11 ou superior
(testado em 3.11, 3.12, 3.13 e 3.14); 4~GB de RAM (pico observado inferior a
1~GB); 3~GB de disco; rede apenas durante a instalação das dependências;
navegador atual apenas para a interface web; usuário comum, sem privilégios de
administrador.

\textbf{Tempo esperado.} Instalação das dependências entre 60 e 120 segundos;
reprodução completa em cerca de 4 minutos no ambiente descrito, e até 12 minutos
em máquina modesta.

\subsection{Dependências}

Todas as dependências de execução são instaladas a partir de
\texttt{requirements-frozen.txt}, que fixa 71 pacotes por versão exata e contém
1.161 hashes SHA-256. A instalação usa \texttt{pip install -{}-require-hashes},
de modo que qualquer divergência de bytes interrompe a instalação em vez de
produzir um ambiente diferente do relatado. O mesmo conjunto está disponível como
\texttt{uv.lock} para o gerenciador \texttt{uv}, que o script de reprodução
prefere quando o encontra.

\textbf{As dependências da interface web estão nesse mesmo arquivo}
(\texttt{streamlit}, \texttt{altair}, \texttt{watchdog}). Não existe passo
adicional de instalação: o script exibe essa informação explicitamente na
primeira etapa. O arquivo \texttt{requirements-web.txt} é apenas uma declaração
legível do subconjunto web e não é usado por nenhum caminho de instalação.

Principais versões fixadas: \texttt{pandas} 3.0.3, \texttt{pyarrow} 24.0.0,
\texttt{pydantic} 2.13.4, \texttt{httpx} 0.28.1, \texttt{beautifulsoup4} 4.15.0,
\texttt{click} 8.4.2, \texttt{rich} 15.0.0, \texttt{pyyaml} 6.0.3,
\texttt{streamlit} 1.58.0, \texttt{altair} 6.2.2, \texttt{pytest} 9.1.1 e
\texttt{ruff} 0.15.20. A lista completa, com hashes, está no repositório.

\textbf{Recursos de terceiros.} Nenhum é necessário para a reprodução. Os
serviços externos (DBLP, Semantic Scholar, OpenAlex, CrossRef, ACM Digital
Library, IEEE Xplore, USENIX e NDSS) são usados apenas pelos comandos opcionais
de coleta, que não fazem parte da reprodução e não alteram o snapshot publicado.
Todos são consultados por endpoints públicos e sem autenticação.

\subsection{Preocupações com segurança}

\textbf{A execução do artefato não oferece risco ao avaliador.} O artefato não
requer privilégios de administrador; não instala serviços, não altera
configurações do sistema e não abre portas externas (a interface web escuta
apenas em \texttt{localhost}); não executa código de terceiros baixado em tempo
de execução, pois as dependências são fixadas por hash; não coleta nem transmite
dados pessoais do avaliador. Após a instalação das dependências, a reprodução é
offline. Todos os arquivos criados ficam dentro do diretório clonado, e remover
esse diretório desfaz completamente a instalação. O snapshot é aberto em modo
somente leitura (\texttt{mode=ro\&immutable=1}), de forma que a execução não pode
corromper o corpus verificado.

\section{Instalação}

\noindent Linux e macOS:
\begin{verbatim}
git clone https://github.com/sidneibarbieri/topvenues-tool
cd topvenues-tool
bash reproduce.sh --profile security-20
\end{verbatim}

\noindent Windows (PowerShell):
\begin{verbatim}
git clone https://github.com/sidneibarbieri/topvenues-tool
cd topvenues-tool
powershell -ExecutionPolicy Bypass -File .\reproduce.ps1 -Profile security-20
\end{verbatim}

\noindent Alternativa por contêiner:
\begin{verbatim}
docker compose up
\end{verbatim}

O script cria um ambiente virtual isolado, instala as dependências fixadas,
materializa o snapshot, verifica seu SHA-256 contra o manifesto, executa a suíte
de testes, renderiza a interface, confere as reivindicações do artigo, compara a
busca ranqueada com a linha de base e exporta uma amostra BibTeX. Ao final desse
comando a ferramenta está instalada e verificada.

\textbf{Observação sobre o Windows.} Se a interface web estiver aberta, o sistema
operacional mantém o arquivo do banco bloqueado e a materialização falha com
\texttt{CorpusBusyError}. Feche o Streamlit e execute o script novamente. A
materialização é protegida por um lock atômico e multiplataforma, de modo que
duas execuções simultâneas se serializam em vez de corromper o banco.

\section{Teste mínimo}

\begin{verbatim}
python -m src.cli --profile security-20 stats
python -m src.cli --profile security-20 search --title "intrusion detection"
python -m src.cli --profile security-20 search --rank "memory corruption" --limit 10
python -m src.cli --profile security-20 export --title intrusion \
    --format bibtex --output amostra.bib
python -m streamlit run web/app.py
\end{verbatim}

\noindent Resultado esperado: o primeiro comando imprime
\texttt{Total Papers: 20305} e \texttt{With Abstracts: 17491 (86.14\%)}; o
segundo exibe 50 registros (limite padrão, de 157 correspondências); o terceiro
retorna registros ordenados por relevância BM25; o quarto grava um arquivo com
2.928 linhas; o quinto abre a interface em \texttt{http://localhost:8501} com
cinco páginas navegáveis. Tempo total aproximado: 90 segundos.

\section{Experimentos}

Todas as reivindicações abaixo são verificadas por um único comando, que também
é executado dentro dos scripts de reprodução:

\begin{verbatim}
python scripts/verify_paper_claims.py --profile security-20
\end{verbatim}

\noindent Recursos comuns a todos os experimentos: menos de 1~GB de RAM, menos
de 3~GB de disco, sem rede e sem GPU.

\subsection{Reivindicação \#1 --- 20.305 artigos entre 2017 e 2026}
Seção 4 do artigo. Tempo: 20~s. Esperado: \texttt{Claim \#1 ... observed 20305} e
\texttt{Claim \#2 ... observed '2017-2026'}.

\subsection{Reivindicação \#2 --- Escopo declarado de 20 veículos}
Seção 3 do artigo. O escopo é dado por \texttt{config.yaml}, fora do código.
Tempo: 20~s. Esperado: \texttt{Claim \#3 ... observed 20}.

\subsection{Reivindicação \#3 --- 17.491 registros com resumo (86,1\%)}
Seção 4 do artigo. Tempo: 20~s. Esperado: \texttt{Claim \#4 ... observed 17491} e
\texttt{Claim \#5 ... observed 86.1}.

\subsection{Reivindicação \#4 --- Todo registro carrega entrada BibTeX}
Seção 5 do artigo. Tempo: 30~s. Esperado: \texttt{Claim \#6 ... observed 20305} e
um arquivo BibTeX com 2.928 linhas.

\subsection{Reivindicação \#5 --- Identidade verificável offline por SHA-256}
Seção 6 do artigo. Comando:
\texttt{python scripts/verify\_profile\_snapshot.py -{}-profile security-20}.
Tempo: 15~s. Esperado: o hash confere com o manifesto
(\texttt{5a35bd6e...7493b08}).

\subsection{Reivindicação \#6 --- Busca ranqueada frente à linha de base}
Seção 5 do artigo. Este experimento traz o \textbf{referencial de comparação}: a
busca por substring (\texttt{LIKE}), equivalente ao que uma planilha ou um
\texttt{grep} fazem, medida contra a busca ranqueada FTS5/BM25 sobre o mesmo
corpus e na mesma máquina. Comando:
\texttt{python scripts/benchmark\_search.py -{}-profile security-20 -{}-trials 11}.
Tempo: 40~s.

\begin{table}[h]
\centering
\begin{tabular}{lrrr}
\hline
Consulta & Linha de base \texttt{LIKE} & Ranqueada BM25 & Ganho \\
\hline
machine learning & 1.962 res., 37,6 ms & 2.093 res., 27--32 ms & $\sim$1,3$\times$ \\
fuzzing & 661 res., 36,0 ms & 661 res., 6,1 ms & $\sim$5,9$\times$ \\
intrusion detection & 337 res., 65,5 ms & 354 res., 3,9 ms & $\sim$16,8$\times$ \\
ransomware & 84 res., 55,8 ms & 84 res., 1,5 ms & $\sim$37$\times$ \\
\hline
\end{tabular}
\end{table}

As contagens são determinísticas e se repetem a cada execução; os tempos são
medianas de 11 repetições e variam entre máquinas. O que o experimento demonstra
é a relação entre as duas colunas.

\subsection{Reivindicação \#7 --- Suíte de testes executada offline}
Seção 6 do artigo. Comando: \texttt{python -m pytest -q}. Tempo: 10~s. O artigo
submetido informa 238 testes, número correto na data da submissão; desde então o
artefato recebeu correções e a suíte cresceu, sem que nenhum teste fosse
removido. As contagens do corpus permanecem idênticas às do artigo, como as
reivindicações \#1 a \#4 demonstram. O número corrente é declarado no
\texttt{README.md} e verificado automaticamente.

\section{Informações adicionais}

\textbf{Documentação.} O \texttt{README.md} do repositório segue o modelo mínimo
do CTA e contém todas as seções exigidas, incluindo solução de problemas para
Windows. A documentação em inglês está em \texttt{README.en.md}.

\textbf{Identificação das reivindicações no artefato.} O script
\texttt{scripts/verify\_paper\_claims.py} associa cada afirmação numérica do
artigo à seção que a enuncia e à consulta que a verifica, e é executado dentro
de \texttt{reproduce.sh} e \texttt{reproduce.ps1}.

\textbf{Demonstração.} Um vídeo de 7min49s em 1920x1080, com narração em inglês
e legendas em português e inglês, acompanha o repositório em
\texttt{docs/assets/demos/} e cobre instalação, verificação offline, busca,
exportações e as análises da interface.

\textbf{Dataset publicado.} A exportação em Parquet do corpus está disponível em
\url{https://huggingface.co/datasets/sidneibarbieri/topvenues}, com o
identificador de perfil, o hash do snapshot e a versão de origem registrados no
cartão do dataset.

\textbf{Licença.} Código sob licença MIT. Metadados bibliográficos e entradas
BibTeX provenientes do DBLP, sob os termos CC0 do DBLP.

\end{document}
